\documentclass[11pt]{book}
\usepackage[paperwidth=145mm,paperheight=200mm,top=18mm,bottom=20mm,inner=20mm,outer=16mm]{geometry}
\usepackage{brumfiel}

% errata.tex -- back-matter list of the 1960 printing's own slips.
%
% Editorial policy (OPEN-QUESTIONS 1.1).  The standing rule for this
% restoration is FAITHFUL TRANSCRIPTION: every slip the compositor made in
% 1960 is set here exactly as the book sets it, with a source comment at its
% site in the tex so no later pass "corrects" it silently.  This appendix is
% the other half of that policy -- it keeps the text faithful while still
% telling a reader who hits a wrong figure reference what the book meant.
%
% Two entries are exceptions, corrected in the text because the transcription
% would otherwise contradict itself on the page; both are called out below.
%
% Every entry was confirmed at 300-450 dpi, most of them twice (teacher's
% photo PDF + Caleb's iPad scan of his own copy).  Entries are keyed to the
% chapter and section, theorem or exercise rather than to a page, since this
% edition is re-typeset and its pagination differs from the original's; the
% original book page is given in parentheses for anyone reading alongside a
% physical copy.

% ---- local macros (hoist into book preamble) ----
\newlist{erratalist}{description}{1}
\setlist[erratalist]{style=nextline, leftmargin=1.4em, labelsep=0.4em,
                     itemsep=4pt, topsep=4pt, parsep=0pt,
                     font=\normalfont\bfseries\small}
% "reads X / meant Y" pair, set compactly under an entry head
\newcommand{\ERprints}[1]{\emph{Prints:} #1}
\newcommand{\ERmeans}[1]{\emph{Read:} #1}
% ---- end local macros ----

\begin{document}
\setcounter{chapter}{18}

\chapter*{\raggedleft\Huge ERRATA}
\addcontentsline{toc}{chapter}{Errata in the 1960 printing}
\markboth{ERRATA IN THE 1960 PRINTING}{}
\vspace{-1em}\noindent\rule{\linewidth}{0.6pt}\vspace{1em}

The text of this edition follows the 1960 printing exactly, including the
places where that printing is itself in error. This list gathers those
places, so that a reader who is sent to the wrong figure, or who meets a
theorem whose statement does not say what its proof needs, can see at once
that the fault is the book's and not the argument's. Page numbers in
parentheses are the original book's.

Two of the entries below are the only ones corrected in the text rather than
reproduced, because in each case the surrounding prose as transcribed would
visibly contradict the page: they are marked \textsc{corrected}.

\begin{erratalist}

\item[Chapter 2 \quad Logic]
\ERprints{Exercise Group 2--12, No.~10 refers to opposite \emph{lines} of the
line $l$ (p.~33).}\\
\ERmeans{opposite \emph{sides}.} The introduction to the same exercise group
also misspells ``equivalence''; that spelling is normalised here.\\
Section 2--8 (p.~35) repeats a clause of the form ``if both of these
statements are true'' twice over --- a compositor's duplication.

\item[Chapter 4 \quad Congruence of segments]
\ERprints{``in \emph{the} Section 4--2'' (p.~64).}\\
\ERmeans{``in Section 4--2''.}

\item[Chapter 5 \quad Measurement of segments]
Theorem 5--4 (p.~95) leaves its third length without the overbar carried by
the first two. The book's use of the segment overbar is erratic throughout,
and is reproduced as printed everywhere; this is simply the most visible
instance.

\item[Chapter 6 \quad Congruence of angles and triangles]
A stray comma stands between ``segments'' and ``are either congruent''
(p.~119). Several sites in this chapter also mix $=$ and $\cong$; that
mixing is the book's own and is reproduced.\\
The self-congruence $\tri{ABE} \cong \tri{ABE}$ on p.~184 is \emph{not} an
error --- it is a genuine dissection argument in the Greek style, and must
not be ``fixed''.

\item[Chapter 7 \quad Use of the congruence theorems]
\ERprints{Construction 7--3 (p.~126) states $ABC \cong A\text{\textup{'}}B\text{\textup{'}}C\text{\textup{'}}$
without the triangle sign on the first name.}\\
\ERmeans{$\tri{ABC} \cong \tri{A'B'C'}$.}\\
Exercise Group 7--8, No.~8 (p.~130) refers to the sides of $A$ where it means
the sides of $\ang{A}$.

\item[Chapter 8 \quad Parallel lines]
\ERprints{The second proof (p.~135) refers the reader to Fig.~8--2.}\\
\ERmeans{Fig.~8--3, which is where angles 1 and 2 are drawn.}\\
\textsc{corrected}: the cross-reference to ``Exercise Group 8--1'' (p.~147)
is set here as \textbf{8--11}. It is a navigation label, and following the
misprint would send a reader to a group twenty pages earlier.

\item[Chapter 9 \quad Similarity of triangles and polygons]
Exercise 9--33 pairs $D$ and $E$ the opposite way round from the exercise
that sets it up; reproduced as printed.

\item[Chapter 10 \quad Area]
\ERprints{A reference to Fig.~9--2 (p.~180).}\\
\ERmeans{Fig.~10--2.} A transcriber's footnote marks this site in the text.

\item[Chapter 11 \quad Circles and regular polygons]
\ERprints{Exercise Group 11--5, No.~7 twice names the centre of the inscribed
\emph{polygon} (p.~197).}\\
\ERmeans{the inscribed \emph{circle}, in both places.}\\
Shanks's computation of $\pi$ is dated to the seventeenth century (p.~204);
it was completed in 1873. The heading ``Remark. 2.'' (p.~205) is normalised
here to ``Remark 2.''

\item[Chapter 13 \quad Loci and sets]
Definition 13--2 (p.~231) defines the union as all points of $S_1$
\emph{and} $S_2$ where it means \emph{or}; the loose phrasing is the book's.\\
Theorem 13--7 (p.~233) names its three parallelograms in inconsistent cyclic
order.\\
Review Exercise 3 (p.~235) bars the second and third $AB$ but not the first.

\item[Chapter 14 \quad Space geometry]
\ERprints{Theorem 14--12 (p.~242) writes the ratio with $A'B$ in the
denominator.}\\
\ERmeans{$A'B'$ --- a prime is missing.}\\
Definition 14--6 (p.~244) names the polygon's last vertex $P_{n-1}$ where it
means $P_n$.\\
The hint to Theorem 14--13 (p.~244) names $\ang{CAV} \cong \ang{BAV}$, both
angles at $A$, where the argument needs the angles at $V$.

\item[Chapter 15 \quad Analytic geometry]
\textsc{corrected}: in Fig.~15--14 (pp.~253\,ff.) the book labels $P_2$ with
the same coordinates as $Q$. It is set here as $(x_2, y_2)$, because the
transcribed prose alongside it refers to $y_2$ explicitly and the figure as
printed contradicts it.\\
Exercise 15--11, No.~1 (p.~262) gives three vertices that do not form a right
triangle. No single-digit change reconstructs one, so this is a genuine error
in the problem rather than a typesetting slip; it is left as printed.\\
Ex.~15--7, No.~16 (p.~257) is \emph{not} an error: the second point reads
$(-5, -7)$, confirmed against the physical copy.

\item[Chapter 16 \quad Philosophy of mathematics]
Berkeley is placed in the seventeenth century (p.~268); he belongs to the
eighteenth. The perturbations of Uranus are dated early in the eighteenth
century (p.~270); the work was done in the 1840s. A footnote on p.~268 names
Baltimore twice in one imprint.

\end{erratalist}

\vspace{1em}
\noindent\emph{Note on footnote marks.} From p.~259 the book's footnotes fall
back to \dag\ because the asterisk is already committed to marking optional
starred material. That is by design, and is preserved.

\end{document}
